% ===========================================================================
\section{Validation}
\label{sec:validation}
% ===========================================================================

There is nothing to validate this sector against directly. No Pakistani facility reports
landfill gas, no drain is monitored for dissolved methane, and no survey measures how much
uncollected waste burns. What exists is other people's models, and comparing a model against a
model establishes consistency rather than accuracy. This section says what that comparison
does and does not show.

\subsection{Against EDGAR}

EDGAR reports Pakistani waste---disposal, burning and wastewater together---at
\SI{26.57}{\mega\tonne} \COe{} for 2023, year-matched to this inventory's reference year
\citep{edgar2025booklet}. This inventory reports \SI{37.45}{\mega\tonne}, or
\SI{141}{\percent} of it. (EDGAR's 2024 value, \SI{27.21}{\mega\tonne}, is carried in the
benchmark file for release-vintage context and is not the comparison basis.)

That row has changed sides during this work, which is what makes it informative. Before
domestic wastewater was built, this sector covered disposal and burning only and read
\SI{49}{\percent} of EDGAR; the gap was then a measure of what the missing category would add.
With 4D1 in, the sector overshoots. The same comparison has therefore diagnosed an omission
and then a possible overestimate, which is more than a single ratio usually earns.

The likeliest single cause is the sewer methane correction factor. At IPCC's stagnant-sewer
\num{0.5} the wastewater category is \SI{24.46}{\mega\tonne} and the sector overshoots; at the
river-discharge \num{0.1} it is \SI{15.39}{\mega\tonne}, the sector is
\SI{28.38}{\mega\tonne}, and the comparison lands at \num{1.07} times EDGAR. No other single
parameter in the sector can move it that far.

It would be easy to treat that as confirmation and adopt \num{0.1}. We do not, and
Section~\ref{sec:discussion} argues why: EDGAR is itself a model making its own assumption
about the same unmeasured quantity, so agreement between the two would demonstrate that two
models had made the same choice rather than that the choice is right.

\subsection{Against Climate TRACE}

Climate TRACE's Pakistani waste subsectors are built from a different evidence base again,
including facility detection. Its solid-waste figures are broadly consistent with the disposal
category reported here; its coverage of open burning and domestic wastewater is not
constructed on a comparable scope, so a total-to-total comparison would mislead rather than
inform and is not reported as one.

\subsection{Against household survey evidence}

The 2024--25 Household Integrated Economic Survey asks every sampled household how its
waste is collected or disposed of \citep{hies2025}, and its design-weighted answers are the
first Pakistani survey evidence bearing on the partition's load-bearing parameters. They
disagree with the What a Waste collection coverage the partition is built on, in one
direction: \SI{42.0}{\percent} of urban households report no collection of any kind against
the \SI{30}{\percent} assumed, and \SI{93.4}{\percent} of rural households against
\SI{70}{\percent}. The route detail supports the open-dumping reading of the uncollected
stream---\SI{88.5}{\percent} of rural households report disposal to open space, and under
\SI{1}{\percent} to waterways.

Two caveats keep this a comparison rather than a correction. Household route shares are not
waste-mass shares: the survey weights households, not tonnes, and converting one to the
other needs stratum generation rates it does not carry. And counting public-bin use as
collected overstates collection where bins are cleared infrequently, which the same module
records. The configuration is therefore priced as a named structural scenario in the
supplement rather than adopted; because the partition conserves generated waste, moving the
uncollected fractions to the survey values moves emissions from slow landfill decay to
immediate burning without changing the waste the country produces.

\subsection{Internal checks that do have teeth}

Three checks are enforced on every run and, unlike the external comparisons, can fail.

\textbf{The mass balance closes.} Disposed plus uncollected equals generated, exactly.
This cannot detect a wrong collection rate, but it can and does detect the class of error
where a change to one category silently fails to appear in the other.

\textbf{The burned quantity is exactly the burned fraction of the uncollected remainder.}
Asserted with equality against an independently recomputed value, so a future refactor that
accidentally applies the burned fraction to total generation rather than to the uncollected
part fails immediately. That specific error would inflate open burning by a factor of two and
is entirely plausible.

\textbf{Provincial allocation reconciles exactly.} The four provinces and the territories sum
to the national figure for each category, not merely for the sector.

\subsection{What would constitute real validation}

Three measurements, in order of what they would settle. Dissolved-methane and flux sampling on
urban drains across a seasonal cycle would resolve the correction factor that dominates the
sector. Landfill gas monitoring at two or three large sites would test the first-order-decay
parameters directly. A waste-characterisation survey at the same sites would replace the
regional composition figures with Pakistani ones and firm up the fossil carbon fraction that
governs reportable \COtwo{} from burning.

None is expensive relative to what the sector contributes to national uncertainty. That
argument is made quantitatively in Section~\ref{sec:discussion}.
